\section{The Last Thing the Composer Does}
\label{sec:ch17-the-last-thing-the-composer-does}
\index{composer!resolvability graph|(}%

Satisfiability is the check chapter~\ref{ch:composition} spent a section on and
this one can finally explain. Before any of that, there is a fact about
\emph{when} it runs that changes how you should read every composition failure
you have ever seen.

\index{composer!satisfiability runs last}%
It runs last, and only if nothing else went wrong. In
\code{federation-factory.ts}, at the end of composition:

\begin{minted}[fontsize=\footnotesize,breaklines]{typescript}
    // Return any composition errors before checking whether all fields are resolvable
    if (this.errors.length > 0) {
      return { errors: this.errors, success: false, warnings: this.warnings };
    }
\end{minted}

The comment is theirs. Every merge-time error, which is to say every
shareability complaint, every type mismatch, every missing or mismatched key
that chapter~\ref{ch:composition} catalogued, is collected and returned before
the resolvability check is even attempted. So a schema pair that fails with one
shareability error has told you nothing at all about whether its fields can be
reached. Fix that error and the next run can produce a satisfiability failure
that was there the whole time. I had read composition output for eight chapters
without knowing that, and it is not something the output can teach you, because
a run that never reached the check looks exactly like a run that passed it.

\subsection*{What the walk is walking}

\index{composer!graph nodes are per subgraph and type}%
The composer builds a graph, once, and threads the same instance through every
subgraph's normalisation pass. Its nodes are keyed by a pair:

\begin{minted}[fontsize=\footnotesize,breaklines]{typescript}
export class Graph {
  edgeId = -1;
  entityDataNodeByTypeName = new Map<TypeName, EntityDataNode>();
  nodeByNodeName = new Map<NodeName, GraphNode>();
  nodesByTypeName = new Map<TypeName, Array<GraphNode>>();
  resolvedRootFieldNodeNames = new Set<NodeName>();
  rootNodeByTypeName = new Map<RootTypeName, RootNode>();
\end{minted}

A node name is \code{subgraph.Type}, so \code{catalog.Product} and
\code{pricing.Product} are two different nodes rather than one merged type.
Ordinary fields are edges inside a subgraph; a resolvable \code{@key} is an edge
between subgraphs. Chapter~\ref{ch:composition} drew that picture already, in
figure~\ref{fig:ch09-resolvability-walk}, and the picture was right. What it did
not know is that the traversal is an ordinary recursive depth-first walk with
cycle detection done by stamping a per-walk index into each edge, rather than
anything cleverer.

\subsection*{Why one route gets the blame}

\index{composer!reports the first failing route}%
Here is chapter~\ref{ch:composition}'s open question. Catalog declares four root
fields that reach a product: \code{products}, \code{browseProducts},
\code{productById} and \code{productBySku}. Make Pricing's key unresolvable and
the composer names \code{Query.products} and nothing else.

The answer is in \code{Graph.validate()}, which loops over root fields and does
not accumulate:

\begin{minted}[fontsize=\footnotesize,breaklines]{typescript}
        const fieldData = getOrThrowError(rootNode.fieldDataByName, rootFieldName, 'fieldDataByName');
        const rootFieldData = newRootFieldData(rootNode.typeName, rootFieldName, fieldData.subgraphNames);
        // If there are no nested entities, then the unresolvable fields must be impossible to resolve.
        if (!involvesEntities) {
          return {
            errors: generateRootResolvabilityErrors({
              unresolvablePaths: rootFieldWalker.unresolvablePaths,
              resDataByPath: rootFieldWalker.resDataByPath,
              rootFieldData,
            }),
            success: false,
          };
        }

        const result = this.validateEntities({ isSharedRootField, rootFieldData, walker: rootFieldWalker });
        if (!result.success) {
          return result;
        }
\end{minted}

Two \code{return} statements, both inside the loop. The first failing root field
ends the function, and \code{validate()} is called from exactly one place and
its result used as-is, so nothing anywhere retries the rest. The loop runs over
a JavaScript \code{Map}, which iterates in insertion order, and edges are
inserted while the parser walks each subgraph's document. Insertion order is
declaration order in the SDL.

\index{composer!the control that moved the blame}%
That is a satisfying explanation and satisfying explanations are exactly the
ones this book has been wrong about before, so I ran the control rather than
printing it. Four compositions, each making the same edit to Pricing, varying
only Catalog's \code{type Query}:

\begin{minted}[fontsize=\footnotesize,breaklines]{text}
all four, as committed                 names ["Query.products"]
products deleted                       names ["Query.browseProducts"]
products and browseProducts deleted    names ["Query.productById"]
all four, productBySku moved to the top names ["Query.productBySku"]
\end{minted}

The last line is the one that does the work. Nothing is removed; one field moves
to the top of the type and the blame moves with it. All four runs report the
same two fields unresolvable, \code{price} and \code{shippingCost}, because the
fault never changed and only the route to it did.

\index{composition errors!fixing the named route surfaces the next}%
What this costs you in practice is worth being blunt about. If you fix the route
the error names, you have not necessarily fixed anything. You have made the
composer walk one field further and find the next way in. On a type with four
entry points that is four rounds of compose, read, edit before the error goes
away, and each round looks like progress.

\subsection*{Two errors, one route, and neither number means what it looks like}

\index{composition errors!repeated identically}%
Chapter~\ref{ch:composition} also recorded that the composer prints some errors
more than once, identically, and established no pattern. There are two different
things going on and they are worth separating, because one of them is not
repetition at all.

The unresolvable key produces two errors, one per unresolvable field, and each
error renders the same explanation. So the sentence naming the root field
appears twice while still naming a single route:

\begin{minted}[fontsize=\footnotesize,breaklines]{text}
unsatisfiable-key: 7 distinct sentences, 10 lines
  the root-field sentence appears x2, naming one route
\end{minted}

The incompatible shared field is the other shape, and there the whole error is
printed twice over, every sentence of it:

\begin{minted}[fontsize=\footnotesize,breaklines]{text}
incompatible-type: 5 distinct sentences, 10 lines
  x2  Each instance of a shared field must resolve identically across ...
  x2  The field "Product.averageRating" could not be federated due to ...
  x2  The discrepancies are as follows:
  x2  The named type "Int" is returned by the following subgraph: "catalog".
  x2  The named type "Float" is returned by the following subgraph: "reviews".
\end{minted}

The two long sentences are shortened by the case rather than by the composer,
which is why they end in an ellipsis.

\index{composer!has no error deduplication}%
Nothing in the composer collapses either. The error list is a bare array, every
call site is a plain \code{push}, and wgc calls composition once and prints one
row per entry, so the doubling starts inside the composer rather than in the
tool around it. I can say that the incompatible-type message is constructed
twice; I cannot say from which two places, and I am not going to guess at it in
print. Both counts are cases in the companion repository, so a release that
starts deduplicating will fail the gate rather than quietly making this section
wrong.

The practical rule is the one chapter~\ref{ch:composition} reached for and can
now be stated with a reason behind it: count faults, never lines. Both cases
above are ten lines of composition output, both are one mistake, and the ten
lines get there by two different routes.

\medskip
That is the composer, which decides what the graph is. The rest of this chapter
is the router, which decides what to do about it, and the first thing it does is
take your document apart.

\index{composer!resolvability graph|)}%
